\title{Large Language Models Can Automatically Engineer Features for Few-Shot Tabular Learning}

\begin{abstract}
Large Language Models (LLMs) possess impressive capacities to address complex and novel reasoning tasks, making them highly promising tools for tabular learning—a foundational component in many practical scenarios. This study introduces a new in-context learning strategy, \textsf{FeatLLM}, which leverages LLMs as feature constructors to generate representations of tabular data that facilitate superior prediction accuracy. Features created by this framework serve as input to straightforward machine learning algorithms, such as linear regression, enabling effective few-shot learning. Crucially, \textsf{FeatLLM} relies solely on this simple predictive model and the transformed features during inference, without requiring real-time interaction with the LLM for every test instance. Compared to prior LLM-centric methods, \textsf{FeatLLM} removes the dependency on sending per-sample queries to an LLM at inference and only necessitates access to the LLM’s API, bypassing typical limitations like prompt length. Through experiments on a variety of tabular datasets from distinct fields, \textsf{FeatLLM} is shown to produce high-quality decision rules, achieving notable improvements (averaging 10\%) over competing frameworks such as TabLLM and STUNT.
\end{abstract}

\section{Introduction}
Large Language Models (LLMs) have recently exhibited remarkable proficiency in generating high-quality solutions for unfamiliar tasks, often without the need for specialized fine-tuning~\cite{creswell2022selection,imani2023mathprompter,taylor2022galactica}. Through training on vast quantities of textual data, LLMs accumulate broad, general knowledge that can supplant the expertise traditionally required from human specialists in a range of domains~\cite{Genomes2021,singhal2023large,wilcox2003role}. This attribute has facilitated progress in numerous automation-related applications, such as discourse analysis~\cite{finch2023leveraging} and program synthesis or correction~\cite{fan2023automated}. In particular, when a prompt contains relevant task instructions alongside a small set of demonstration cases, LLMs are able to grasp contextual cues and discern significant variables~\cite{brown2020language,zhao2023survey}, reflecting their aptitude for advanced data interpretation—akin to the function of a domain expert in feature selection and construction.

Harnessing LLMs’ broad generalist reasoning has proven fruitful within tabular learning settings. For example, domain insights—like associating advanced age with increased disease vulnerability—can enhance the performance of predictive models for medical diagnosis. Recent methodologies in this arena express both feature and task requirements in natural language, transform tabular data into serialized input formats, and employ these structured prompts in LLM pipelines~\cite{dinh2022lift,wang2023anypredict}. These strategies are frequently coupled with in-context learning paradigms~\cite{nam2023semi} or lightweight adaptation techniques such as parameter-efficient fine-tuning~\cite{dinh2022lift,hegselmann2023tabllm}. The resulting integration of LLM-derived prior knowledge has yielded improved predictive outcomes over classical tabular learners, especially in data-scarce (“low-shot”) situations~\cite{hegselmann2023tabllm}.

Existing approaches for tabular learning utilizing LLMs face several notable drawbacks. First, performing end-to-end predictions necessitates at least one LLM inference for each individual sample, leading to considerable computational overhead. Second, most strategies depend on fine-tuning the LLM to attain strong performance, which is impractical for LLMs lacking accessible training tools or tuning interfaces—a significant issue given that state-of-the-art LLMs increasingly restrict user access to inference-only APIs~\cite{anil2023palm,openai2023gpt4}. Third, the majority of current techniques struggle with lengthy prompts~\cite{liu2023lost}; as the feature count in tabular datasets rises, resulting in extended serialized textual inputs, many methods become infeasible or experience degraded accuracy. Collectively, these obstacles substantially impede the practical deployment of LLM-driven tabular learning models.

In contrast, our method advocates for a new paradigm: rather than relying on LLMs for direct end-to-end prediction, we harness their strengths in feature engineering. By proposing the creation of more informative features that incorporate LLMs’ prior knowledge, we present a novel in-context learning methodology called \textsf{FeatLLM}. This framework seeks to extract the latent “criteria” embedded in LLM decision-making. As an illustration, for disease prediction, the LLM can infer and output explicit decision rules describing which feature combinations indicate the presence of a disease. These criteria can be used to build new features, which effectively replace conventional features in downstream learning tasks. Consequently, the LLM’s role transitions to generating high-level features, facilitating a streamlined downstream process—when robust features are available, simple models suffice, thereby reducing inference time. The inherent in-context learning abilities of LLMs enable us to generate these features simply by adding demonstration examples to the prompts, without any additional model training. Moreover, employing ensemble-based strategies such as feature bagging~\cite{breiman2001random} allows us to circumvent problems arising from large prompt sizes.

\textsf{FeatLLM} organizes the prompt for feature engineering into two main components: (1) analyzing the task to identify the connections between features and targets, and (2) leveraging both prior insights and limited exemplar data to construct rule-based criteria for class differentiation. By following this sequential reasoning strategy, the LLMs exclude irrelevant features during rule generation. The resulting rules are implemented on the dataset (interpreted as program instructions), converting the original samples into binary indicators reflecting rule adherence. Subsequently, a simple model is trained to predict class probabilities using these newly generated features. Model stability is enhanced through ensemble techniques, with repeated feature creation and feature bagging. Adjusting the quantity of features and samples during bagging helps control prompt size, preventing scalability issues. Because \textsf{FeatLLM} operates without conventional training requirements and maintains low inference overhead, it is applicable to a wide range of practical scenarios.

Our assessment of \textsf{FeatLLM} encompasses 13 distinct tabular datasets in settings with few training samples, demonstrating its reliability and effectiveness. Results indicate that LLMs adeptly incorporate both background knowledge and information extracted from the dataset, yielding high-caliber features. The proposed methodology achieves superior performance compared to current few-shot learning methods across multiple benchmarks. Code is available at the anonymized GitHub repository:~\texttt{https://github.com/Sungwon-Han/FeatLLM}.

\section{Related Work}

\subsection{Few-Shot Learning with Tabular Data} The progress in few-shot learning methodologies has made it possible to derive robust and transferable representations from datasets with very limited labeled examples~\cite{chen2018closer}. Although early approaches prioritized image data, the capacity of few-shot learning to generalize has extended across multiple modalities, such as text, audio, and more recently, tabular datasets~\cite{majumder2022few,nam2023stunt,schick2021exploiting}. Relying on only a handful of labeled examples may result in models picking up on spurious associations~\cite{bartlett2021deep}. To mitigate these challenges, contemporary techniques integrate supplementary information or utilize domain-specific prior knowledge to induce beneficial inductive biases during model development. Approaches include leveraging extensive unlabeled data collections with carefully designed augmentations in semi-supervised settings~\cite{ucar2021subtab,yoon2020vime} and employing unsupervised, task-independent pre-training to optimize model initialization~\cite{somepalli2021saint}. Unsupervised pre-training objectives often entail procedures like partial masking or data corruption, combined with reconstruction loss functions~\cite{arik2021tabnet,majmundar2022met}, or adopting contrastive objectives through diverse augmentation strategies~\cite{bahri2022scarf,chen2020simple}. However, because tabular data exhibits substantial heterogeneity, it is difficult to devise universally suitable augmentation methods. As a result, such techniques have generally shown limited utility in the few-shot context~\cite{nam2023stunt}.

Emerging research advocates training models on an expansive and varied set of tabular datasets, which may be entirely unrelated to the specific target task, then using transfer learning to adapt to new tasks~\cite{zhang2023generative,levin2022transfer,zhu2023xtab}. For example, TransTab~\cite{wang2022transtab} implements transformer-based architectures that learn to capture semantic associations across columns by utilizing column names in addition to cell values. The representations learned from this wide array of tables and columns allow for zero-shot and few-shot predictions for unseen tabular tasks. Similarly, TabPFN~\cite{hollmann2022tabpfn} synthesizes data via combinations of diverse distributions, which is then used for pre-training a Bayesian neural network. Post-pre-training, the model conducts inference using both the available training and test data for the target task without further finetuning. The accumulated prior knowledge from multiple datasets has enabled TabPFN to outperform standard tree-based techniques and demonstrated its advantage in settings with very limited supervision.

\subsection{Language-Interfaced Tabular Learning} Integrating large language models (LLMs)~\cite{anil2023palm,openai2023gpt4} provides another avenue for encoding and exploiting prior knowledge. Trained on vast and wide-ranging textual datasets, LLMs have shown remarkable adaptability, generalizing effectively to unfamiliar tasks and evidencing utility across numerous fields~\cite{brown2020language}. Recent works have sought to tap into the extensive prior knowledge inherent in LLMs for tabular data tasks. These methods commonly involve transforming tabular information into textual representations and then incorporating these serializations within the prompts supplied to LLMs~\cite{dinh2022lift,hegselmann2023tabllm,wang2023anypredict}. By embedding data, descriptive features, and task context into these prompts, LLMs can infer task-feature relationships needed for predictions. Advances in this area encompass both parameter-efficient tuning mechanisms, such as LoRA~\cite{hu2021lora} and IA3~\cite{liu2022few}, as well as in-context learning strategies, where few-shot demonstration examples are appended to input prompts in lieu of explicit model training~\cite{dinh2022lift,hegselmann2023tabllm,nam2023semi,slack2023tablet}.

Although the prior techniques demonstrate notable success in enhancing few-shot tabular learning, reliance on running inference exclusively through the LLM presents obstacles for adoption in industrial contexts. This work seeks to move away from the traditional end-to-end black-box paradigm, where every sample is processed through the LLM for prediction. Rather, the proposed approach leverages the LLM to infer class-specific conditions or rules. \textsf{FeatLLM} converts these rules into programmatic code, generating novel features and facilitating predictions independent of the LLM. Additionally, it exploits in-context learning for rule extraction by leveraging previous knowledge and information derived from the few-shot examples.

\section{Methods}
\textsf{FeatLLM} focuses on identifying ``rules"—i.e., conditions linked with each class label—using the limited examples provided, as illustrated in Figure~\ref{fig:main_model}. These rules are parsed by the LLM, which then transforms them into binary feature indicators for each sample, reflecting whether or not each rule is satisfied. The collection of these binary features is multiplied by a set of non-negative, trainable weights via a dot product operation to estimate the probability for each class. By optimizing this model, the relative influence of each feature on class probability is determined (refer to Section~\ref{Sec:training}). The procedure is repeated with bagging, and predictions are combined using ensemble techniques. The formal problem setup and detailed descriptions of each stage are provided in the following subsections.

\vspace*{-0.12in}
\paragraph{Problem formulation.} We analyze a tabular dataset containing $N$ labeled entries, denoted as $\mathcal{D} = \{(\mathbf{x}^i, \mathbf{y}^i)\}_{i=1}^{N}$. Each $\mathbf{x}^i$ is a vector with $d$ input features, and each feature is associated with a natural language name, represented as $F = \{f_j\}_{j=1}^{d}$, accompanied by independent definitions. In this classification context, each label $\mathbf{y}^i$ belongs to a predetermined set of classes. For $k$-shot learning scenarios, the model is trained using only $k$ randomly chosen labeled samples, where $k<N$.

\subsection{Prompt Design for Extracting Rules}
\label{Sec:prompt_design}
To foster accurate rule extraction from the LLM, we structure the prompt to imitate the expert approach taken by humans in tabular analysis. The prompt consists of three principal elements (illustrated in Fig.~\ref{fig:prompt_for_rule} and detailed in Appendix~\ref{sec:example_rule}):

\vspace*{-0.12in}
\paragraph{Basic information description.} This section delivers the fundamental data required for effective problem-solving (highlighted in orange in Fig.~\ref{fig:prompt_for_rule}). It contains a task outline (including label details), feature explanations, and exemplar demonstrations. The task outline is posed as a direct question (for instance, ``Does this patient's myocardial infarction complication information suggest chronic heart failure? Yes or no?"). Additional examples appear in Appendix~\ref{sec:task_desc}. Feature explanations specify the value type and may optionally add definitions. For categorical features, examples of potential categories are included. Demonstrative examples entail serializing select training instances into text paired with their true labels. The serialization format follows approaches used in previous research~\cite{dinh2022lift,hegselmann2023tabllm}:

\begin{align}
&\text{Serialize}(\mathbf{x}^i,\mathbf{y}^i, F) = \nonumber \\ 
&``\ f_1\ \text{is}\ \mathbf{x}^i_1.\ ... \ f_d\  \text{is}\  \mathbf{x}^i_d.\ \ \text{Answer:}\  \mathbf{y}^i\ ”
\end{align}

\vspace*{-0.12in}
\paragraph{Reasoning instruction.} Our objective is to strengthen the reasoning capabilities of LLMs by supplying targeted reasoning instructions (as illustrated in the blue text in Fig.~\ref{fig:prompt_for_rule}). The instruction commences with a sentence akin to that used in the chain-of-thought methodology~\cite{wei2022chain}. Subsequently, the reasoning workflow is segmented into two phases. Initially, the LLM is prompted to deduce causal relationships or tendencies between features and the specified task description, utilizing only general knowledge or common sense—without reference to demonstration examples. This step maximizes the LLM’s ability to leverage existing knowledge about the task. In the second phase, example demonstrations together with the knowledge distilled in the initial phase are used by the LLM to induce specific rules for each class. By segmenting reasoning in this manner, we help the model steer clear of irrelevant correlations and enhance its focus on pertinent features. To control for underfitting or unnecessarily verbose prompts due to fluctuating rule counts, we limit the extraction to a predetermined number of rules (10 per class)\footnote{A discussion regarding rule count is provided in Section~\ref{sec:analysis}.}. Furthermore, during rule generation, the LLM is instructed to consult the type information of each feature as described in the initial information section, thereby mitigating the risk of mismatched feature types in the resulting rules. 

\vspace*{-0.12in}
\paragraph{Response instruction.} To streamline rule extraction and subsequent application, explicit instructions are given to the LLM on formatting its responses (see yellow text in Fig.~\ref{fig:prompt_for_rule}). The prompt specifies guidelines for organizing responses by answer class (i.e., response structure), and for the form each rule should take (i.e., rule format). This guidance empowers the LLM to select suitable formats according to each feature’s type. Absent further constraints, the LLM is also enabled to integrate multiple rules using logical connectives such as AND and OR. Illustrative results can be found in Appendix~\ref{sec:outcome-rule}.

\subsection{Inferring Class Likelihood via Rules}
\label{Sec:training}

\paragraph{Rule parsing for feature construction.} In this phase, we transform the rules generated earlier into binary features. For every class, we assign these features to represent whether a sample meets the criteria defined by the rules for that class (yielding either ‘0’ or ‘1’). These features are valuable for estimating the probability that a sample pertains to a specific class. Since the LLM’s output consists of natural language, an automated strategy for parsing this output and constructing features is necessary. Although we restrict the format of responses and rules to streamline parsing, outputs from the LLM may still incorporate unpredictable syntactic variations~\cite{kovavc2023large}. For example, the model may modify comparison operators, replacing symbols like “$>$” or “$<$” with equivalent words such as “greater than” or “less than,” thereby complicating the parsing process.

To overcome difficulties posed by noisy textual data, rather than resorting to elaborate code implementations, we harness the capabilities of the LLM itself. The LLM excels at grasping semantic information regardless of syntactic inconsistencies and can generate succinct code from directives, owing to its training on programming data. To capitalize on these strengths, we prepare prompts comprising function names, input/output descriptions, and the relevant rules, then submit them to the LLM. Illustrations in Figures~\ref{fig:prompt_for_parsing} and ~\ref{sec:example_parsing}-\ref{sec:example_parse_outcome} (see Appendix) display examples of such prompts along with their results. The resulting code snippets are executed by invoking Python’s \texttt{exec()} function to facilitate data processing.

\vspace*{-0.12in}
\paragraph{Estimating class probability.} After obtaining new binary features derived from class-specific rules, one rudimentary way to assess the likelihood that a sample belongs to a given class is to sum the number of rules it meets for each class (i.e., accumulate the values of the binary features within each class). Yet, the influence of rules and features varies, making it necessary to quantify their relative importance by learning from the training data. We address this by fitting a standard machine learning model to the binary features corresponding to each class, thereby estimating their contribution. As an example, we might employ a bias-free linear model. Let the binary feature vector created from sample $\mathbf{x}^i$ for class $k$ be $\mathbf{z}_k^i \in \{0, 1\}^R$, where $R$ is the total number of rules for each class, and $c$ is the number of classes overall. We then determine the probability of each class $\mathbf{p}^i$ via:

\begin{align}
\text{logit}_k^i &= \max(\mathbf{w}_k, 0) \cdot \mathbf{z}_k^i, \nonumber \\
\mathbf{p}^i &= \text{Softmax}({[\text{logit}_{1}^i, ..., \text{logit}_{c}^i]}). \label{eq:infer_prob}
\end{align}

In this framework, $\mathbf{w}_k$ denotes the adjustable parameters within the linear model and reflects the relevance assigned to each input feature. By constraining these weights to a positive domain, the approach guarantees that features extracted by the LLM will not detract from a class's probability during the training phase. Subsequently, the logit values are transformed into probability distributions over classes utilizing the softmax mechanism. The model parameters $\mathbf{w}_k$ are refined via cross-entropy loss, with optimization performed on a limited number of few-shot examples.

\vspace*{-0.12in}
\paragraph{Ensembling with bagging.} To conclude, the methodology involves iteratively repeating the described process, resulting in several inference models whose outputs are aggregated to generate a final ensemble prediction. The ensemble leverages the set of trained weights from $T$ independent runs, written as $\{\mathbf{w}[t]\}_{t=1}^T$. The prediction for each class $\mathbf{p}^i$ is obtained by averaging the probabilities calculated with each $\mathbf{w}[t]$, according to Eq.~\ref{eq:infer_prob}.

For the purpose of diversifying rule generation during ensemble construction, a high temperature is selected during LLM inference. Additionally, the arrangement of few-shot examples is shuffled, recognizing that permutation in demonstration sequence can affect LLM response outcomes~\cite{lu2022fantastically}\footnote{Further examination of rule diversity can be found in Appendix~\ref{sec:diversity}}. To maximize predictive variation for increased resilience, bagging is employed, which involves sampling subsets of features or data instances for each trial—a strategy that becomes increasingly effective with larger training sets or feature numbers. The ensemble strategy proposed here offers two primary benefits. Firstly, erroneous rules produced in certain trials by the LLM may be offset by correct rules from other iterations, thereby bolstering the framework's reliability. This concept is related to the self-consistency property of LLMs~\cite{wang2022self}: rules frequently produced across different trials have a higher probability of being valid. Secondly, the ensemble approach mitigates restrictions associated with LLM prompt length; when tabular datasets exceed allowable prompt capacities, bagging serves to curtail input size, helping to maintain model robustness. While any individual rule set may not encompass all feature relationships, the ensemble of multiple weak learners derived from repeated trials counteracts gaps in feature and instance representation across separate models.

\section{Experiments}
We assess \textsf{FeatLLM} across a range of tabular datasets, undertaking several analyses to explore the mechanisms and efficacy of our proposed framework. Due to space limitations, supplementary experiments—such as exploring alternative LLMs, qualitative evaluations, and additional analyses—are presented in the Appendix.

\subsection{Performance Evaluation}
\paragraph{Datasets.}
Our study encompasses 13 datasets designed for binary and multi-class classification tasks: (1) Adult~\cite{asuncion2007uci}, used to forecast whether an individual’s annual income surpasses \$50,000; (2) Bank~\cite{moro2014data}, aiming to determine the likelihood of a customer subscribing to a term deposit; (3) Blood~\cite{yeh2009knowledge}, which predicts repeat blood donation behavior; (4) Car~\cite{kadra2021well}, for assessing car quality; (5) Communities~\cite{misc_communities_and_crime_183}, employed to estimate crime rates within specific areas; (6) Credit-g~\cite{kadra2021well}, targeting credit risk evaluation; (7) Diabetes\footnote{\texttt{kaggle.com/datasets/uciml/pima-indians-diabetes-database}}, identifying diabetic cases; (8) Heart\footnote{\texttt{kaggle.com/datasets/fedesoriano/heart-failure-prediction}}, predicting coronary artery disease occurrences; and (9) Myocardial~\cite{misc_myocardial_infarction_complications_579}, which identifies chronic heart failure presence.

Additionally, we integrate both contemporary and synthetic datasets rarely seen in prior literature and absent from LLM pre-training sets: (10) Cultivars, which predicts the grain yield classification of soybean cultivars\footnote{\texttt{archive.ics.uci.edu/dataset/913}}; (11) NHANES, used to categorize individuals’ age groups based on records\footnote{\texttt{archive.ics.uci.edu/dataset/887}}; (12) Sequence-type, a task distinguishing sequences among arithmetic, geometric, Fibonacci, and Collatz types; and (13) Solution-mix, determining whether the concentration level in a mixture of four solutions exceeds 0.5. The first two datasets were deposited to the UCI Machine Learning Repository after September 2023, affirming their absence from LLM API (gpt-3.5-turbo-0613) training data utilized in our experiments. The final two are synthetic datasets, precluding memorization by the LLM API and allowing unbiased evaluation.

The datasets feature substantial diversity in sample counts and feature dimensionality, as summarized in Table~\ref{Tab:basic-desc}. Every dataset comes with explicit attribute labels and descriptions. Datasets such as ‘Communities’ and ‘Myocardial’ each comprise over 100 features, presenting distinct hurdles linked to LLMs’ restricted context window.

\vspace*{-0.12in}
\paragraph{Baselines.}
Nine baseline methods are incorporated for comparison. The initial three are standard tabular classifiers: (1) Logistic regression (LogReg), (2) XGBoost~\cite{chen2016xgboost}, and (3) RandomForest~\cite{ho1995random}. The subsequent two operate on unlabeled data: (4) SCARF~\cite{bahri2022scarf}, and (5) STUNT~\cite{nam2023stunt}; note that acquiring substantial unlabeled data may not always be practical in real-world contexts. (6) TabPFN~\cite{hollmann2022tabpfn} is a recent approach leveraging synthetic datasets with varied distributions for model pretraining. The final three are based on LLMs: (7) In-context learning (In-context)~\cite{wei2022emergent}, (8) TABLET~\cite{slack2023tablet}, and (9) TabLLM~\cite{hegselmann2023tabllm}. In-context learning incorporates few-shot samples directly within prompts without model modification. TABLET enhances prompts by introducing supplementary hints through external classifiers with rule-based sets and prototypes, aiming for improved inference quality. TabLLM leverages parameter-efficient adaptation methods, such as IA3~\cite{liu2022few}, for few-shot tuning of LLMs.

\vspace*{-0.12in}
\paragraph{Implementation details.}
\textsf{FeatLLM} leverages GPT-3.5 as its primary LLM, although its architecture remains flexible and compatible with other LLM variants (see Appendix~\ref{sec:backbones} for performance across different backbones). During inference, the LLM operates with a temperature parameter of 0.5, while the top-p is set to the API default of 1. The method incorporates 20 ensembles and extracts 10 rules by default. Further insights into the sensitivity to hyper-parameters can be found in Figure~\ref{fig:hyperparameter2} and Appendix~\ref{sec:hyper-parameter}.
For training the linear model, Adam is employed with a learning rate of 0.01, and training is conducted over 200 epochs. To select the optimal number of training epochs, $k$-fold cross-validation is applied. When reproducing baseline methods, GPT-3.5 serves as the model for in-context learning, while TabLLM utilizes T0~\cite{sanh2022multitask} consistent with its original configuration. It should be noted that TabLLM is limited to LLMs with accessible checkpoints, preventing the use of GPT-3.5 in this setting. As for conventional machine learning models, including LogReg, XGBoost, and RandomForest, hyper-parameter tuning is achieved using Grid-search coupled with $k$-fold cross-validation. The $k$ value is fixed at either 2 or 4 to ensure every class is represented in the training set. Additional implementation information can be found in Appendix~\ref{sec:baselines}.

\vspace*{-0.12in}
\paragraph{Results.}
Tables~\ref{tab:main1} and \ref{tab:main2} provide comparative results across several datasets. Each experiment was run three times using varied random splits of training and test data, reporting both the mean and standard deviation for the resulting AUC metrics. In most cases, \textsf{FeatLLM} either achieved the highest performance or placed second among competing models. Remarkably, our method matched the accuracy of traditional tabular learning baselines using only 4-shot samples, while those baselines required 32 shots to achieve similar results (as depicted in Fig.~\ref{fig:compares}). Although methods such as STUNT and TabPFN achieved notable gains for selected datasets, their overall advantages over traditional machine learning techniques were marginal. Furthermore, approaches leveraging LLMs—including in-context learning, TABLET, and TabLLM—were unable to process tabular datasets featuring a high dimensionality of features. In contrast, our architecture, which integrates bagging and ensembling strategies, remained robust and delivered strong results even when confronted with high-dimensional feature spaces. It is also worth highlighting that our ensemble method, in principle, could be extended to other LLM-based algorithms. However, applying this adaptation would result in doubling the inference computations per sample, thus significantly amplifying computational demands to impractical levels.

\subsection{Analysis \& Discussion}
\label{sec:analysis}
\paragraph{Ablation study.} We investigate the effect of principal framework components through ablation experiments, focusing on: (1) \textsf{-Tuning}: removing the weight adjustment step for the linear model and, instead, summing the binary feature values directly to derive the logit; (2) \textsf{-Ensemble}: excluding the ensemble method entirely; (3) \textsf{-Description}: omitting explicit feature descriptions; and (4) \textsf{-Reasoning}: bypassing the first step in the reasoning instruction process when rule generation takes place.

The outcomes of these ablations are summarized in Table~\ref{tab:ablation}, which reports average changes in AUC values across datasets. Disabling any key module led to reduced predictive performance. Of note, the weight tuning procedure demonstrated an increasing benefit as the shot number grew, implying that with a more extensive training set, rule importance can be estimated more precisely, amplifying the gains from tuning. Conversely, the inclusion of feature descriptions and reasoning instructions proved more critical under low-shot conditions, underscoring the role of leveraging LLM's prior knowledge in scarce data regimes. Ultimately, the ensemble component reliably produced performance improvements irrespective of the experimental setup.

\vspace*{-0.12in}
\paragraph{Training \& inference time comparison.} We evaluate and contrast the time required for both training and inference across various methods. All tests are carried out using the Adult dataset with a training set comprising 16 examples. For most methods, computations are run on a single A100 GPU; however, TabLLM utilizes four A100 GPUs for model parallelism. Table~\ref{tab:cost} lists the comparative runtimes. \textsf{FeatLLM} stands out by delivering relatively fast inference speeds, similar to conventional baselines. On the other hand, LLM-driven approaches such as In-context learning, TABLET, and TabLLM exhibit significantly slower inference phases. Regarding training durations, \textsf{FeatLLM} performs on par with alternative LLM-based techniques, requiring only 30 API calls. This minimal API usage does not strain computational or financial resources and can be further expedited with parallel execution.

\vspace*{-0.12in}
\paragraph{Quality of features generated by \textsf{FeatLLM}.}
We perform an analysis to assess how informative the rule-based features produced by \textsf{FeatLLM} are for real-world tasks. The approach entails calculating the average AUROC value between the features created and the target labels, followed by determining the proportion of features considered useful—those yielding an AUROC greater than 0.5—across datasets. Table~\ref{tab:quality} reports these findings. The data show that most generated rules are pertinent to the target variable and contribute meaningful information for task resolution (see the ``Before tuning'' column). Furthermore, when tuning the linear model using dataset-specific information, rules with low association (those whose weights fall below a predefined cutoff) are systematically eliminated (see ``After tuning''). The threshold, set at 0.1, was determined by examining the distribution of feature weights.

\vspace*{-0.12in}
\paragraph{Effect of adding spurious correlations.} To evaluate how well different approaches remove noisy, spurious correlations in situations with limited labeled data, we performed a targeted experiment. This experiment leveraged the Heart dataset, supplemented by random columns drawn from the Adult dataset, which bear no relevance to the Heart prediction task. The analysis tracked how the models' accuracy shifted as the quantity of extraneous columns was systematically increased. For consistency, all methods were given exactly 8 labeled samples and were not provided with any unlabeled data; as a result, algorithms such as SCARF or STUNT were not included.

Figure~\ref{fig:spurious} illustrates the impact of spurious correlations on model performance. Among the tested methods, \textsf{FeatLLM} exhibited minimal reduction in accuracy attributable to these irrelevant features. This resilience can be traced to the framework’s reliance on prior knowledge to determine which features are pertinent when generating rule sets, thereby minimizing the creation of rules based on unrelated columns and enhancing robustness. In practice, rules utilizing noisy columns were found to comprise just 1--2\% of the overall rule pool. However, it is important to note that \textsf{FeatLLM} is not immune to the pitfalls of learning spurious associations or confusing causal relations if random columns are sourced from datasets with genuine task similarity (see Appendix~\ref{sec:spurious-appendix}).

\vspace*{-0.12in}
\paragraph{Hyper-parameter analysis}
\textsf{FeatLLM} incorporates two primary hyper-parameters: the ensemble size and the quantity of rules extracted per class in each LLM invocation. To evaluate how these parameters affect model outcomes, we ran a series of tests. Our observations reveal that increasing the ensemble size generally enhances predictive performance, yet this improvement plateaus near the value of 20 (refer to Fig.~\ref{fig:appendix-hyperparameter1} in the Appendix). Regarding the rule count, no statistically significant variations were identified across different values, though selecting 10 strikes an optimal balance between prompt length and accuracy (as illustrated in Fig.~\ref{fig:hyperparameter2}). We hypothesize that augmenting the number of rules escalates input dimensionality, thus expanding the information available, but simultaneously increases the likelihood of overfitting when data is limited.

\section{Conclusion}
We introduce \textsf{FeatLLM}, which establishes new benchmarks in few-shot tabular learning using LLMs. Rather than leveraging LLMs solely for instance-level predictions, \textsf{FeatLLM} centers on generating predictive criteria, emulating the function of feature engineers. Through the combined mechanisms of rule induction and ensemble bagging, \textsf{FeatLLM} achieves low inference overhead, necessitates only API access and not model training, and surmounts challenges associated with extensive feature dimensionality. The method delivers robust predictive capability and offers an advantageous solution for deploying LLMs efficiently on practical tabular datasets.

\textsf{FeatLLM} is tailored for scenarios with limited data availability. In subsequent work, we aim to broaden its applicability by facilitating automated feature engineering for datasets with greater sample sizes, exploring other feature extraction methods beyond rule-based systems, and integrating interpretability to expand its reach across diverse analytical tasks.

\section{Broader Impact} The proposed framework, \textsf{FeatLLM}, leverages LLMs’ embedded prior knowledge to elevate performance on tabular learning tasks, surpassing traditional supervised methods in data efficiency. By supplementing sparse training data with domain knowledge, this approach offers a promising pathway to reducing overfitting risks. Its utility is especially pronounced in applied fields such as finance and healthcare, where labeling poses considerable practical challenges, and pretrained LLMs contain substantial relevant information. For broader adoption, a key area for future research will involve assessing the influence of societal biases and misinformation present within LLMs’ prior knowledge, as this information can significantly shape predictive outcomes and may overshadow empirical labeled data.

\end{document}